\section{Late labels and the delayed loop}\label{sec:loop}

\subsection{The soft sensor and its interval}\label{sec:loop:interval}

A soft sensor predicts a quality variable $y$ from process measurements $x$. At every calibrator
update $t$, spaced $\Delta u$ apart, it issues a prediction $\hat y_t$; the laboratory or the
analyzer supplies $y_t$ later. The residual is $r_t = y_t - \hat y_t$, and its absolute value, the
score $s_t = |r_t|$, measures how wrong the sensor was. The interval issued at update $t$ is
centred on the prediction,
\begin{equation}
  C_t = [\,\hat y_t - q_t,\; \hat y_t + q_t\,],
  \label{eq:interval}
\end{equation}
with the half-width $q_t$ equal to the $\lceil (1-\alpha_t)(W+1)\rceil$-th smallest of the $W$
most recent scores whose labels have arrived, the split-conformal quantile
\citep{papadopoulos2002}. The half-width is infinite when that rank exceeds $W$, that is for
$\alpha_t < 1/(W+1)$, and zero for $\alpha_t \ge 1$. The level $\alpha_t$ adapts online to hold
a target miscoverage $\alpha^\ast$: a miss lowers it, which widens later intervals, and a hit
raises it \citep{gibbs2021aci}. Throughout, $\alpha^\ast = 0.10$ and $W = 200$ (40 labelled
events in the arm that updates only on new analyzer values; Section~\ref{sec:study}).

\subsection{Dead time}\label{sec:loop:deadtime}

Let the label of a sample become available a time $D$ after the interval for that sample was
issued. With updates every $\Delta u$, the label is $a = \lceil D/\Delta u \rceil$ updates old when
it is first read. Labels are read at the start of an update, before the interval is issued, so the
label of the interval issued at update $s$ first moves the level used at update $s+\theta+1$, with
the dead time
\begin{equation}
  \theta = a - 1 .
  \label{eq:theta}
\end{equation}
A calibrator that issued before reading would have $\theta = a$. With $\theta = 0$ the scheme is
ordinary adaptive conformal inference; \citet{elhalabi2026} write $\tau = \theta + 1$. The level
then obeys
\begin{equation}
  \alpha_{t+1} = \alpha_t + \gamma\,(\alpha^\ast - e_{t-\theta}),
  \qquad e_s = \mathbf{1}\{s_s > q_s\},
  \label{eq:loop}
\end{equation}
with step size $\gamma > 0$, and the score $s_{t-\theta}$ enters the window when its label
arrives. The miss $e_s$ is judged against the half-width $q_s$ issued for sample $s$, which the
calibrator stores until the label arrives: stored pairing. Judging it against the half-width in
force when the label arrives, arrival pairing, is a bookkeeping error that an implementation
holding a single last half-width makes easily; the study uses it as a negative control
(Section~\ref{sec:study}). Fig.~\ref{fig:block} shows the loop.

\begin{figure}[t]
\centering
\begin{tikzpicture}[
    font=\scriptsize, >={Stealth[length=1.5mm]}, semithick,
    blk/.style={draw, rounded corners=1pt, align=center, inner sep=2.5pt,
      minimum height=7mm},
    sum/.style={draw, circle, inner sep=0pt, minimum size=3.6mm}]
  \node[blk] (model) at (0,0) {point\\model};
  \node[sum] (sum) at (1.25,0) {$+$};
  \node[blk] (int) at (2.9,0) {interval\\$\hat y_t + \hat r_t \pm q_t$};
  \node[blk] (comp) at (1.25,-1.4) {compensator\\$(\theta{+}1)$ ahead};
  \node[blk] (quant) at (2.9,-1.4) {quantile\\$1{-}\alpha_t$ of\\the window};
  \node[blk] (miss) at (4.85,-1.4) {miss $e_s$,\\score $s_s$,\\residual $r_s$};
  \node[blk] (lab) at (6.5,-1.4) {laboratory\\$z^{-\theta}$};
  \node[blk] (integ) at (2.9,-2.85) {integrator, gain $\gamma$\\$\alpha \leftarrow \alpha + \gamma(\alpha^\ast - e_s)$};
  \draw[->] (-1.0,0) node[left] {$x_t$} -- (model);
  \draw[->] (model) -- node[above] {$\hat y_t$} (sum);
  \draw[->] (sum) -- (int);
  \draw[->] (int.east) -- (5.75,0) node[right] {$C_t$};
  \draw[->] (4.85,0) -- node[left, align=right] {stored\\$C_s$} (miss.north);
  \draw[->] (6.5,0.45) node[above] {sample} -- (lab.north);
  \draw[->] (lab) -- node[above] {$y_s$} (miss);
  \draw[->] (miss) -- node[above] {$s_s$} (quant);
  \draw[->] (quant) -- node[right] {$q_t$} (int);
  \draw[->] (integ) -- node[right] {$\alpha_t$} (quant);
  \draw[->] ([xshift=-2.5mm]miss.south) |- node[pos=0.3, right] {$e_s$} (integ.east);
  \draw[->] ([xshift=3mm]miss.south) -- ++(0,-1.95) -| node[pos=0.25, below] {$r_s$, $s \le t-\theta-1$} (comp.south);
  \draw[->] (comp) -- node[left] {$\hat r_t$} (sum);
  \fill (4.85,0) circle (1.1pt);
\end{tikzpicture}
\caption{The calibrated soft sensor with delayed labels. The interval is centred on the point
prediction plus the predicted residual $\hat r_t$ ($\hat r_t = 0$ without compensation,
Section~\ref{sec:remedies:comp}); its half-width is the conformal quantile at level
$1 - \alpha_t$ of the last $W$ scores. The label $y_s$ arrives $\theta$ updates after the interval
for sample $s$ was issued and moves the level one update later, through an integrator of gain
$\gamma$; its miss $e_s$ is judged against the stored interval $C_s$. The calibrator is an
integral loop with dead time $\theta$.}
\label{fig:block}
\end{figure}

\subsection{Documented delays}\label{sec:loop:delays}

Table~\ref{tab:delays} derives $\theta$ from documented delays for the two processes of the
study. In the debutanizer, the butane content of the bottoms is measured by a gas chromatograph
on the overhead of the downstream column, with a 15-min cycle and about 45~min from sample to
result \citep[Sec.~6.2]{fortuna2007}; the same book also gives 20 to 60~min
\citep[Sec.~5.5]{fortuna2007}, and 30 to 75~min is reported by \citet[p.~42]{siddharth2019}. The
book models the seven inputs of the public data set at 12-min updates (Sec.~6.2) but analyses a
recording of the same column at 6-min updates (Sec.~4.1), with outlier counts out of 2394 rows,
the public file's length; it does not say which applies, so $\theta = 3$ is primary and
$\theta = 7$ secondary. The Tennessee Eastman product analyzer has a sampling interval and a dead
time of 0.25~h \citep[Table~5]{downs1993}. The same delay gives a dead time of 4 at 3-min
updates and of 14 at 1-min updates, and of 0 if the calibrator updates only when a new analyzer
value appears: the faster the calibrator runs on held labels, the longer its dead time.

\begin{table}[t]
\centering\footnotesize
\caption{Dead time from documented delays: label age $a = \lceil D/\Delta u \rceil$ and
$\theta = a - 1$ (Eq.~\eqref{eq:theta}). Debutanizer: gas chromatograph with a 15-min cycle,
about 45~min \citep[Sec.~6.2]{fortuna2007}; the range combines 20 to 60~min
\citep[Sec.~5.5]{fortuna2007} and 30 to 75~min \citep[p.~42]{siddharth2019}; the book gives
12-min and 6-min updates without saying which fits the public data set. Tennessee Eastman:
product analyzer XMEAS(40) \citep[Table~5]{downs1993}; $D$ counts from the row whose inputs give
the prediction to the row in which the label is written (counted from the composition's capture,
a reading is one row older when it first appears). Primary cells in bold.}
\label{tab:delays}
\begin{tabular}{@{}llrrr@{}}
\toprule
Process & Delay $D$ & $\Delta u$ & $a$ & $\theta$ \\
\midrule
Debutanizer & about 45 min & 12 min & 4 & \textbf{3} \\
            & about 45 min & 6 min & 8 & 7 \\
            & 20 to 75 min & 12 or 6 min & 2 to 13 & 1 to 12 \\
Tennessee Eastman & 0.25 h & 3 min & 5 & \textbf{4} \\
                  & 0.25 h & 1 min & 15 & 14 \\
                  & 0.25 h & per new value & 1 & 0 \\
\bottomrule
\end{tabular}
\end{table}

\subsection{What is measured}\label{sec:loop:metrics}

The primary measure is the uninformative share $U$: the share of issued intervals that are
infinite or wider than the range of $y$ over the calibration data. An operator cannot act on
such an interval, whatever the coverage figure says. Marginal coverage, the share of labelled
intervals that contain $y$, is reported with its Wilson 95\,\% interval \citep{wilson1927};
widths by the median, the 90th and the 99th percentiles of the finite widths; local behaviour by
the worst rolling coverage, the minimum coverage over 100 consecutive intervals. The persistence
of the misses is measured by the lag-1 autocorrelation of the miss indicator $e$, the mean length
of a run of consecutive misses, and the integrated autocorrelation time (IACT),
$1 + 2\sum_{k \ge 1} \rho_e(k)$, summed over an automatic window (Section~\ref{sec:study}). The
IACT is the factor by which serial correlation inflates the variance of a mean of misses, and so
of the level that Eq.~\eqref{eq:loop} integrates from them.
